% !PRINTED 74
\weberpage{74}

Da der untere Index \(\nu\) in \(A_\nu^{(\mu)}\) gar nicht vorkommt, so ändert sich \(A_\nu^{(\mu)}\) nicht, wenn der untere Index \(\nu\) durch einen anderen ersetzt wird. Dann aber verschwindet nach \S~21, VI. die Determinante. Wir erhalten demnach aus (2) die folgende wichtige Relation, in der \(\mu,\nu\) irgend zwei von einander verschiedene Ziffern \(1,2,\ldots,n\) sein können:
\begin{equation}
0=a_\mu^{(1)}A_\nu^{(1)}+a_\mu^{(2)}A_\nu^{(2)}+\cdots+a_\mu^{(n)}A_\nu^{(n)},
\tag{6}
\end{equation}
und ebenso bekommt man aus (3):
\begin{equation}
0=a_1^{(\mu)}A_1^{(\nu)}+a_2^{(\mu)}A_2^{(\nu)}+\cdots+a_n^{(\mu)}A_n^{(\nu)}.
\tag{7}
\end{equation}

Beispielsweise ergiebt sich aus (5), wenn \(a,b,c\) durch \(a',b',c'\) ersetzt werden:
\begin{equation}
a'(b'c''-c'b'')+b'(c'a''-a'c'')+c'(a'b''-b'a'')=0,
\tag{8}
\end{equation}
eine Formel, von deren Richtigkeit man sich durch die einfachste Rechnung überzeugt.

Wenn wir die Relation (6) mit einem beliebigen Factor \(\lambda\) multipliciren und zu (2) addiren, so erhalten wir die Formel:
\begin{equation}
\begin{aligned}
A={}&(a_\nu^{(1)}+\lambda a_\mu^{(1)})A_\nu^{(1)}
+(a_\nu^{(2)}+\lambda a_\mu^{(2)})A_\nu^{(2)}\\
&+\cdots+(a_\nu^{(n)}+\lambda a_\mu^{(n)})A_\nu^{(n)},
\end{aligned}
\tag{9}
\end{equation}
die uns den folgenden Satz ausdrückt:

\gesperrt{VII. Die Determinante ändert ihren Werth nicht, wenn man zu den Elementen einer Zeile, die mit einem beliebigen gemeinschaftlichen Factor multiplicirten entsprechenden Elemente einer anderen Zeile addirt.}

Derselbe Satz gilt auch von den Colonnen. Er wird zur Vereinfachung und numerischen Berechnung von Determinanten oft mit Nutzen verwendet. Wir fügen noch folgende Sätze bei, die sich aus den Darstellungen (2), (3) sofort ablesen lassen.

\gesperrt{VIII. Wenn alle Elemente einer Zeile oder einer Colonne einen gemeinschaftlichen Factor haben, so kann dieser weggelassen und als Factor vor die Determinante gesetzt werden.}

Denn es ist nach (2):
\[
p a_\nu^{(1)}A_\nu^{(1)}+p a_\nu^{(2)}A_\nu^{(2)}+\cdots+p a_\nu^{(n)}A_\nu^{(n)}=pA.
\]

\gesperrt{IX. Wenn in einer Zeile oder in einer Colonne alle Elemente bis auf eines verschwinden, so reducirt}
%CONT>>
